\chapter{Route contract and hypothesis ledger}
\label{route-ch-contract}

\section{Fixed objects and functors}

\begin{definition}[The AJCR curve and its degree-zero Picard functor]
  \label{route-conv-objects}
  \lean{AlgebraicGeometry.genus, AlgebraicGeometry.pic0TypeFunctor}
  \leanok
  \dcref{custom}

  Let \(K\) be a field and let
  \[
    C\longrightarrow \Spec K
  \]
  be smooth of relative dimension one, proper, and geometrically irreducible.
  Put
  \(g=g(C)=\dim_K H^1(C,\mathcal O_C)\).  For a \(K\)-scheme \(T\), write
  \[
    \operatorname{Pic}^0_C(T)=\operatorname{pic0TypeFunctor}(C)(T)
  \]
  for the degree-zero subgroup of the big-etale-sheafified relative Picard
  functor used in Lean.  Its elements are classes whose degree after every
  field-valued pullback is zero.  This notation does \emph{not} assert that
  this functor has already been identified with the identity component of a
  Picard scheme.

  No rational point of \(C\) is assumed.  In particular, no step may introduce
  a global \(K\)-point merely to normalize a divisor or a Picard class.
\end{definition}

\begin{proof}
  This is the exact binder of the public Lean interfaces: \([\mathrm{Field}\ K]\),
  \(C:\Over(\Spec K)\), and the three displayed geometric typeclasses.  The
  degree-zero functor is defined before any representer and is therefore the
  object which every later \(\mathrm{RepresentableBy}\) certificate must
  represent.
\end{proof}

\section{The one permitted route}

The dependency order is fixed as follows.  Every arrow is a mandatory
dependency, not a status claim.  The boxed break marks the first missing
mathematical producer; later arrows describe blocked consumers, not a change
of route.

\[
\begin{gathered}
  (R1)\ D_g\text{ represents the widened degree-}g\text{ divisor functor}\\
  \Downarrow\\
  (R2)\ W\simeq \operatorname{Pic}^{g,\mathrm{rk}=1}_C\\
  \Downarrow\\
  (R3)\ \{W_a\to\operatorname{Pic}^0_C\}_{a}\text{ covers over }\Omega\\
  \Downarrow\\
  (R4)\ J_\Omega\text{ represents }\operatorname{Pic}^0_{C_\Omega}
    \text{ and is qc, qs, lfp}\\
  \Downarrow\\
  (R5)\ P_0=(E_0,\mathrm{GlueData},J_{E_0}),\qquad
       J_{E_0}\times_{E_0}\Omega\simeq J_\Omega\\
  \Downarrow\\[-2pt]
  \boxed{(R6)\ \text{jointly choose }P\text{ with }L=P.N.1/K
    \text{ finite Galois, descend }\xi_L,
    \text{ and prove }h_{P.\mathrm{gluedOver}}
      \simeq\operatorname{Pic}^0_{C_L}}\\
  \Downarrow\\
  (R7)\ \text{prove }P.\mathrm{gluedMap}\text{ projective, hence
       finite-in-affine for this same }J_L=P.\mathrm{gluedOver}\\
  \Downarrow\\
  (R8)\ J=J_L/\operatorname{Gal}(L/K)\text{ represents }
        \operatorname{Pic}^0_C\\
  \Downarrow\\
  (R9)\ \mathrm{PicRepDatum}\longrightarrow\mathrm{JacobianData}
       \longrightarrow\text{Phase 8 consumers}.
\end{gathered}
\]

The route does not pass through a quotient of the unrestricted
admissible-degree Abel map, a symmetric-power/Weil construction, an Albanese
first construction, or a filtered-colimit theorem.  Those branches are
quarantined in Appendix~\ref{route-app-quarantine}.

\section{Hypothesis ledger}

Fix a separable closure \(\Omega/K\): thus \(\Omega/K\) is algebraic and
separable and \(\Omega\) is separably closed.  The exact Lean context includes
\(\operatorname{Algebra.IsSeparable} K\ \Omega\).  Every finite coefficient
set used in descent is required first to lie in a finite subextension
\(E/K\subset\Omega/K\) carrying both \(\operatorname{Module.Finite}_K E\) and
\(\operatorname{Algebra.IsSeparable} K E\).  Only after that certificate is
available may the finite normal-closure theorem be applied.  A stage carrying
only \(\operatorname{Algebra.IsAlgebraic} K E\) is not admissible for R6.
The checked carrier comparison may first be built at an arbitrary finite
subextension \(E_0/K\).  The missing universal-class producer must enlarge all
finitely many coefficients to a finite normal subextension and only then
finalize the legacy package \(P\).
Accordingly its literal final field \(L=P.N.1\) is already finite Galois over
\(K\); no later step changes the carrier.  The following ledger is exhaustive.

\begin{center}
\begin{tabular}{p{0.08\textwidth}p{0.37\textwidth}p{0.43\textwidth}}
\textbf{ID} & \textbf{Hypothesis} & \textbf{Scope and discharge}\\ \hline
H1 & \(K\) is a field. & Global; never discharged.\\
H2 & \(C/K\) is smooth of relative dimension one, proper, and geometrically
irreducible. & Global challenge input; preserved by field base change.\\
H3 & \(\Omega/K\) is algebraic and separable, and \(\Omega\) is separably
closed; every finite stage \(E/K\subset\Omega\) used by the route carries
\(\operatorname{Module.Finite}_K E\) and \(\operatorname{Algebra.IsSeparable} K E\). &
Used for translated rank-one coverage, finite-stage spreading, and the
finite-Galois final stage.  The finite normal closure is taken only after the
separability certificate is present, and remains inside \(\Omega\).  This is
stronger than merely choosing an algebraic separably closed extension over an
imperfect field; the exact Lean context must include
\texttt{Algebra.IsSeparable K Omega}.\\
H4 & A legacy finite-stage glue package \(P_0\) has final field
\(E_0=P_0.N.1\subset\Omega\). & Checked from finite presentation of the finite
atlas; it supplies the comparison shape but is not yet the final route carrier.\\
H5 & A legacy package \(P\) is selected jointly with the universal-class
descent, with \(L=P.N.1\subset\Omega\) finite Galois over \(K\). & The
normal-closure theorem supplies the field enlargement; reflecting the whole
finite diagram at that stage and producing the exact Lean package are planned
parts of the first producer.\\
H6 & \(h_{J_L}\simeq\operatorname{Pic}^0_{C_L}\) on the literal carrier
\(J_L=P.\mathrm{gluedOver}\). & \textbf{Unresolved first mathematical
producer}; it must come from a class \(\xi_L\) on this carrier and its Yoneda
equivalence.\\
H7 & \(P.\mathrm{gluedMap}\) is projective. & \textbf{Unresolved exact-carrier
geometry}; this is the one selected sufficient input for orbit-affineness.\\
H8 & \(J_L\to\Spec L\) is locally of finite type and quasi-compact. & Produced
from the finite affine presentation; the conditional Galois wrappers consume
these certificates.\\
H9 & A \(K\)-rational point \(P\in C(K)\). & Used only by pointed Abel--Jacobi
consumers in Phase 8; never used to construct \(\operatorname{Pic}^0\).\\
H10 & Properness, smoothness of relative dimension \(g\), and geometric
irreducibility of the final carrier. & Outputs still to be instantiated from
the original-field datum; they are not hypotheses of representability.\\
\end{tabular}
\end{center}

\begin{remark}[Status discipline]
  \label{route-rem-status}
  \dcref{custom}

  \emph{Proved} means a public-root-reachable Lean declaration with the
  displayed type.  \emph{Imported} means a source or Mathlib theorem used with
  its full hypotheses.  \emph{Conditional} means Lean checks an implication
  whose antecedent remains visible; a \texttt{\textbackslash leanok} marker on
  such a node certifies only that generic implication, not its unresolved route
  instance.  \emph{Planned} means the exact target type is fixed but no
  producer exists.  \emph{Unresolved} means that target is the first unmet
  mathematical obligation on the route.  A source declaration not reachable
  from the public root is \emph{unverified} here and receives no Lean marker.
\end{remark}

\begin{proof}
  This convention separates implementation state from mathematical truth and
  prevents a consumer theorem from being counted as its missing producer.
\end{proof}
